% =====================================================================
%  Shared sections (model, stability, control, swing-up) used by build_full_report.py and build_phase_reports.py
%  (layout follows the senior reference report: title page, contents,
%   numbered sections, code listings, tabulation, conclusion, contributions)
% =====================================================================
\documentclass[12pt,a4paper]{article}
\usepackage[a4paper,margin=1in]{geometry}
\usepackage[T1]{fontenc}
\usepackage{lmodern}
\usepackage{amsmath,amssymb,bm}
\usepackage{graphicx,float,booktabs,array,multirow}
\usepackage{xcolor}
\usepackage{algorithm}
\usepackage{algpseudocode}
\usepackage{fancyhdr}
\usepackage{enumitem}
\usepackage{caption}
\usepackage[hidelinks]{hyperref}

\pagestyle{fancy}
\fancyhf{}
\renewcommand{\headrulewidth}{0.6pt}
\fancyfoot[C]{\thepage}
\setlength{\headheight}{14.5pt}
\setlength{\parskip}{4pt}
\setlength{\parindent}{0pt}
\newcommand{\Nm}{\,\mathrm{N\,m}}
\captionsetup{font=small}

\begin{document}

% ------------------------------------------------------------------ title
\begin{titlepage}
\centering
\vspace*{2cm}
\includegraphics[width=4.2cm]{figures/logo-000.jpg}\\[3cm]
{\Huge Systems Thinking Project}\\[0.5cm]
{\Large Acrobot: Stability, Control and Swing-Up}\\[0.25cm]
{\large Phases 2, 3 and 4}\\[2.2cm]
\begin{tabular}{cc}
\textbf{Author:} Bhargav & \textbf{Author:} Sreekar\\
ID: 2025102061 & ID: 2025102074\\[0.8cm]
\textbf{Author:} Kaushik & \textbf{Author:} Divijat\\
ID: 2025102071 & ID: 2025102040
\end{tabular}
\vfill
{\large November 11, 2026}
\end{titlepage}

\tableofcontents
\newpage

% =====================================================================
\section{Introduction}
% =====================================================================
\subsection{What is an Acrobot?}
An acrobot is a two-link pendulum that hangs from a fixed pivot. The first link (the ``shoulder'' link) is free to rotate about the pivot; the second link is joined to the end of the first at the ``elbow''. The only motor sits at the elbow, so the shoulder is \emph{passive}. It is the robot version of a gymnast on a high bar who can only swing the legs at the hips. Because there are two degrees of freedom but only one actuator, the acrobot is a standard example of an \emph{underactuated} system [1], [2].

\subsection{Model}
The configuration is described by two generalised coordinates: $q_1$, the angle of link~1 from the downward vertical, and $q_2$, the angle of link~2 relative to link~1. The full state contains the angles and their rates:
\[
\mathbf q=\begin{bmatrix}q_1\\q_2\end{bmatrix},\qquad
\mathbf x=\begin{bmatrix}q_1\\q_2\\\dot q_1\\\dot q_2\end{bmatrix}.
\]
\begin{figure}[H]
\centering
\includegraphics[width=0.42\textwidth]{figures/acrobot_diagram.pdf}
\caption{The acrobot: passive shoulder at the pivot, motor torque $\tau$ at the elbow; dots mark the centres of mass.}
\label{fig:acrobot}
\end{figure}

\subsection{Equations of Motion}
The dynamics (given in the Phase~1 helper and verified in Phase~1) are
\begin{equation}
M(q)\,\ddot q+C(q,\dot q)\,\dot q=\tau_g(q)+B\,\tau,\qquad B=\begin{bmatrix}0\\1\end{bmatrix},\qquad |\tau|\le 20\Nm .\label{eq:eom}
\end{equation}

\textbf{$\tau$ (Input Torque):} the torque of the elbow motor, the only quantity we command. $B=[0,\,1]^T$ says it acts on the elbow equation only.

\textbf{$M(q)$ (Inertia Matrix):} the resistance to angular acceleration. It depends on the elbow angle because the mass distribution changes as the arm folds:
\begin{equation}
M(q)=\begin{bmatrix}2.663+2\cos q_2 & 1.33+\cos q_2\\ 1.33+\cos q_2 & 1.33\end{bmatrix},\qquad \det M=1.7729-\cos^2q_2\ge0.7729>0 .
\end{equation}

\textbf{$C(q,\dot q)$ (Coriolis Matrix):} the coupling forces between the links (centrifugal and Coriolis effects):
\begin{equation}
C(q,\dot q)=\sin q_2\begin{bmatrix}-2\dot q_2 & -\dot q_2\\ \dot q_1 & 0\end{bmatrix}.
\end{equation}

\textbf{$\tau_g(q)$ (Gravity Torque):} the torque that gravity exerts on each joint, $\tau_g=-\partial U/\partial q$, with potential energy $U$:
\begin{equation}
U=-14.715\cos q_1-9.81\cos(q_1+q_2),\qquad
\tau_g(q)=-9.81\begin{bmatrix}1.5\sin q_1+\sin(q_1+q_2)\\ \sin(q_1+q_2)\end{bmatrix}.
\end{equation}

\textbf{System Parameters:}
\begin{gather*}
m_1=m_2=1\,\mathrm{kg},\qquad l_1=1\,\mathrm m,\qquad l_2=2\,\mathrm m,\qquad l_{c1}=0.5\,\mathrm m,\qquad l_{c2}=1\,\mathrm m,\\
I_1=0.083\,\mathrm{kg\,m^2},\qquad I_2=0.33\,\mathrm{kg\,m^2},\qquad g=9.81\,\mathrm{m\,s^{-2}}.
\end{gather*}

\subsection{Equilibria and Linear Model (from Phase 1)}
With $\tau^*=0$ the equilibria satisfy $\tau_g(q^*)=0$, i.e.\ $\sin q_1^*=0$ and $\sin(q_1^*+q_2^*)=0$. This gives four postures:
\begin{table}[H]
\centering
\caption{The four equilibria}
\begin{tabular}{|c|c|c|c|l|}
\hline
 & Posture & $(q_1^*,q_2^*)$ & $U$ [J] & Description\\ \hline
$E_1$ & down--down & $(0,0)$ & $-24.525$ & both links hanging\\
$E_2$ & down--up & $(0,\pi)$ & $-4.905$ & link 1 down, link 2 up\\
$E_3$ & up--up & $(\pi,0)$ & $+24.525$ & fully upright (target)\\
$E_4$ & up--down & $(\pi,\pi)$ & $+4.905$ & link 1 up, link 2 down\\ \hline
\end{tabular}
\end{table}
Linearizing about each equilibrium with $\tilde x=x-x^*$ gives
\begin{equation}
\dot{\tilde x}=A\tilde x+B_\ell\tilde\tau,\qquad A=\begin{bmatrix}0&I\\P&0\end{bmatrix},\qquad B_\ell=\begin{bmatrix}0\\r\end{bmatrix},\label{eq:lin}
\end{equation}
with $P=M(q^*)^{-1}K(q^*)$ and $r=M(q^*)^{-1}B$,
where $K=\partial\tau_g/\partial q$. The transfer functions from $\tau$ to the angles are
\begin{gather}
G_1(s)=\frac{r_1s^2+(P_{12}r_2-P_{22}r_1)}{\Delta(s)},\qquad
G_2(s)=\frac{r_2s^2+(P_{21}r_1-P_{11}r_2)}{\Delta(s)},\label{eq:tf}\\
\Delta(s)=s^4-\operatorname{tr}(P)\,s^2+\det P .\nonumber
\end{gather}
\begin{table}[H]
\centering
\caption{Transfer functions at the four equilibria (Phase 1)}
\label{tab:tf}
\small
\begin{tabular}{|c|l|l|l|}
\hline
 & $\Delta(s)$ & $N_1(s)$ ($\tau\to q_1$) & $N_2(s)$ ($\tau\to q_2$)\\ \hline
$E_1$ & $s^4+42.2410\,s^2+186.7719$ & $-3.0147\,(s^2+4.2103)$ & $6.0332\,(s^2+5.2595)$\\
$E_2$ & $s^4+8.4025\,s^2-186.7719$ & $-0.4270\,(s^2-29.7273)$ & $0.8578\,(s^2+7.3982)$\\
$E_3$ & $s^4-42.2410\,s^2+186.7719$ & $-3.0147\,(s^2-4.2103)$ & $6.0332\,(s^2-5.2595)$\\
$E_4$ & $s^4-8.4025\,s^2-186.7719$ & $-0.4270\,(s^2+29.7273)$ & $0.8578\,(s^2-7.3982)$\\ \hline
\end{tabular}
\end{table}

\subsection{Objective}
\begin{enumerate}
\item \textbf{Phase 2:} determine the stability of every equilibrium using both the transfer functions and the state-space model, support each conclusion with simulation, and decide whether feedback from a \emph{single} measured joint angle can stabilize the upright equilibrium $E_3$.
\item \textbf{Phase 3:} design controllers that stabilize $E_3$ and at least one other equilibrium, justify the design choices, and characterize how far from the equilibrium each controller still works with $\tau_{\max}=20\Nm$.
\item \textbf{Phase 4:} starting from hanging ($E_1$), bring the acrobot to $E_3$ and hold it, and explain with evidence why the controller behaves as it does.
\end{enumerate}

% =====================================================================
\section{Phase 2: Stability Analysis}
% =====================================================================
\subsection{Stability Criteria}
\begin{itemize}
\item \textbf{State space:} the linear model~\eqref{eq:lin} is \emph{asymptotically stable} if every eigenvalue of $A$ has $\operatorname{Re}\lambda<0$, \emph{unstable} if any has $\operatorname{Re}\lambda>0$, and \emph{marginally stable} if the eigenvalues on the imaginary axis are simple and none lies to the right.
\item \textbf{Transfer function:} $G_j$ is \emph{BIBO stable} if every pole has $\operatorname{Re}p<0$.
\item \textbf{Nonlinear system (Lyapunov's indirect method):} an eigenvalue with $\operatorname{Re}\lambda>0$ makes the equilibrium of~\eqref{eq:eom} unstable; eigenvalues on the imaginary axis leave the question open [4].
\end{itemize}

\subsection{Poles of the Transfer Functions = Eigenvalues of $A$}
At every equilibrium
\[
\operatorname{rank}[B_\ell\ \ AB_\ell\ \ A^2B_\ell\ \ A^3B_\ell]=4,\qquad
\operatorname{rank}\begin{bmatrix}C_j\\C_jA\\C_jA^2\\C_jA^3\end{bmatrix}=4\quad(j=1,2),
\]
so both $G_1$ and $G_2$ are minimal realizations with no pole--zero cancellation. Their poles are exactly the eigenvalues of $A$, and the transfer-function and state-space verdicts must coincide.

\subsection{Characteristic Polynomial}
\begin{equation}
\Delta(s)=s^4+0\cdot s^3-\operatorname{tr}(P)\,s^2+0\cdot s+\det P .
\end{equation}
A Hurwitz polynomial has all coefficients positive. Here the $s^3$ and $s^1$ coefficients are zero because the model has no friction, so \textbf{no equilibrium is asymptotically stable with $\tau=0$}. Writing $\lambda=s^2$ (the eigenvalues of $P$):
\begin{equation}
\lambda^2-\operatorname{tr}(P)\lambda+\det P=0,\qquad \lambda<0\Rightarrow s=\pm j\sqrt{-\lambda},\qquad \lambda>0\Rightarrow s=\pm\sqrt\lambda .
\end{equation}
\begin{itemize}
\item $\det P>0$, $\operatorname{tr}P<0$: both $\lambda<0$, all poles on the $j\omega$-axis ($E_1$).
\item $\det P<0$: one $\lambda$ of each sign, one real RHP pole ($E_2$, $E_4$).
\item $\det P>0$, $\operatorname{tr}P>0$: both $\lambda>0$, two RHP poles ($E_3$).
\end{itemize}

\subsection{Link to the Potential Energy}
Since $K=-\nabla^2U$ and $M\succ0$, $P=M^{-1}K$ is similar to $M^{-1/2}KM^{-1/2}$, which is congruent to $K$. By Sylvester's law of inertia,
\begin{equation}
\#\{\text{RHP poles}\}=\#\{\text{negative eigenvalues of }\nabla^2U\}.
\end{equation}
A minimum of $U$ gives no RHP pole, a saddle one, and a maximum two [3].

\subsection{Stability of Each Equilibrium}
\begin{table}[H]
\centering
\caption{Stability of the four equilibria}
\label{tab:stab}
\small
\resizebox{\textwidth}{!}{\begin{tabular}{|c|c|c|c|c|c|}
\hline
 & $(\operatorname{tr}P,\det P)$ & poles of $G_1,G_2$ $=$ eig$(A)$ & eig$\nabla^2U$ & RHP & Verdict\\ \hline
$E_1$ & $(-42.24,\ 186.77)$ & $\pm j2.2400,\ \pm j6.1011$ & $4.91,\ 29.43$ (min) & 0 & marginally stable\\
$E_2$ & $(-8.40,\ -186.77)$ & $\pm3.1775,\ \pm j4.3010$ & $-14.72,\ 9.81$ (saddle) & 1 & unstable\\
$E_3$ & $(42.24,\ 186.77)$ & $\pm2.2400,\ \pm6.1011$ & $-29.43,\ -4.91$ (max) & 2 & unstable\\
$E_4$ & $(8.40,\ -186.77)$ & $\pm4.3010,\ \pm j3.1775$ & $-9.81,\ 14.72$ (saddle) & 1 & unstable\\ \hline
\end{tabular}}
\end{table}

\textbf{$E_1$ (down--down).}
\begin{itemize}
\item \textbf{Transfer function:} the poles are simple and on the imaginary axis, so $G_1$ and $G_2$ are \emph{not} BIBO stable. A bounded input at a resonance, $\tau=\sin(2.24t)$, gives a double pole at $\pm j2.24$ and a response that grows like $t\sin(2.24t)$.
\item \textbf{State space:} the free response is a sum of two undamped modes ($\omega=2.24$ and $6.10$~rad/s): bounded but not decaying, so \textbf{marginally stable}.
\item \textbf{Nonlinear:} linearization is inconclusive, so we use energy as a Lyapunov function:
\begin{equation}
V(x)=\tfrac12\dot q^TM(q)\dot q+U(q)-U(q^*),\qquad \nabla^2U(q^*)\succ0\Rightarrow V\succ0,\qquad \dot V=\dot q^TB\tau=0 .
\end{equation}
$E_1$ is \textbf{stable in the sense of Lyapunov but not asymptotically stable}: the swing never decays.
\end{itemize}

\textbf{$E_2$ and $E_4$ (one link up).} $\det P=-186.77<0$, so each has one real RHP pole ($+3.1775$ at $E_2$, $+4.3010$ at $E_4$). Both are BIBO unstable, $A$ is unstable, and by Lyapunov's indirect method the nonlinear equilibria are unstable.

\textbf{$E_3$ (up--up).} Gravity flips, $P_{E3}=-P_{E1}$, so $\lambda$ changes sign and the poles become $\pm2.2400$, $\pm6.1011$: \textbf{two RHP poles}, unstable in every sense.

\subsection{Simulation Evidence}
All Phase~2--4 simulations use the full nonlinear model in \textbf{PyBullet} [5] (uniform rods, no damping, step $0.25$~ms). Checked against an RK4 integration of~\eqref{eq:eom} with $q(0)=(0.8,-0.5)$~rad and $\tau=3\sin2t$, the two agree to within $0.52^\circ$ over 5~s.
\begin{figure}[H]
\centering
\includegraphics[width=\textwidth]{figures/fig_openloop.pdf}
\caption{Free response ($\tau=0$) near each equilibrium: PyBullet (colour) vs.\ linear model (black).}
\label{fig:ol}
\end{figure}
\begin{figure}[H]
\centering
\includegraphics[width=\textwidth]{figures/fig_growth.pdf}
\caption{(a) Growth of a tiny offset on a log scale; dotted lines have the slope of the largest RHP pole. (b) Spectrum of $\tilde q_1$ near $E_1$; dotted lines are the poles.}
\label{fig:growth}
\end{figure}
\begin{table}[H]
\centering
\caption{Linear predictions vs.\ PyBullet measurements}
\label{tab:evidence}
\small
\begin{tabular}{|c|l|c|c|c|}
\hline
 & Quantity & Predicted & Measured & Error\\ \hline
$E_1$ & oscillation frequencies & $2.240,\ 6.101$ rad/s & $2.239,\ 5.997$ rad/s & within 1 bin\\
$E_1$ & energy change over 60 s & $0$ & $1.2\times10^{-3}$ J & integration error\\
$E_2$ & growth rate & $3.1775$ s$^{-1}$ & $3.171$ s$^{-1}$ & $0.22\%$\\
$E_3$ & growth rate & $6.1011$ s$^{-1}$ & $6.088$ s$^{-1}$ & $0.21\%$\\
$E_4$ & growth rate & $4.3010$ s$^{-1}$ & $4.294$ s$^{-1}$ & $0.15\%$\\ \hline
\end{tabular}
\end{table}
\textbf{Observations:}
\begin{itemize}
\item $E_1$ oscillates forever at its initial amplitude, with frequencies matching the poles.
\item A tiny offset at $E_2$, $E_3$, $E_4$ grows exponentially at exactly the rate of the largest RHP pole. At $E_3$ a $0.01^\circ$ tilt reaches $30^\circ$ in $1.55$~s.
\item The linear model tracks the nonlinear one until about $20^\circ$, then over-predicts.
\end{itemize}

\subsection{Can One Measured Angle Stabilize the Upright?}
\subsubsection{Static or PD Feedback}
Measuring one angle $y=\tilde q_j$ and applying $\tau=-(K_p+K_ds)\,y$ with $G_j=(b_2s^2+b_0)/(s^4+a_2s^2+a_0)$ gives the closed-loop polynomial
\begin{equation}
s^4+(K_db_2)\,s^3+(a_2+K_pb_2)\,s^2+(K_db_0)\,s+(a_0+K_pb_0) .\label{eq:pd}
\end{equation}
Hurwitz requires $K_db_2>0$ \emph{and} $K_db_0>0$, i.e.\ $b_2$ and $b_0$ of the same sign. At $E_3$:
\begin{itemize}
\item measuring $q_1$: $b_2=-3.0147$, $b_0=+12.6926$, opposite signs;
\item measuring $q_2$: $b_2=+6.0332$, $b_0=-31.7316$, opposite signs.
\end{itemize}
The signs are opposite because the zeros $\pm2.0519$ and $\pm2.2934$ are \textbf{real}: one of each pair lies in the right half-plane. \textbf{No $K_p$, $K_d$ can stabilize $E_3$ from one angle.} A grid search confirms it: the smallest attainable largest real part is $0.96$ ($q_1$) and $2.29$ ($q_2$), and in PyBullet even these gains fall from $1^\circ$ within 2~s.
\begin{figure}[H]
\centering
\includegraphics[width=\textwidth]{figures/fig_pd.pdf}
\caption{PD from one angle at $E_3$: (a, b) largest closed-loop real part over the gain plane, positive everywhere; (c) PyBullet with the best gains.}
\label{fig:pd}
\end{figure}

\subsubsection{Dynamic (Observer-Based) Feedback}
Each angle alone makes $(A,C_j)$ observable, so an observer can reconstruct the full state and LQR can act on the estimate:
\begin{equation}
\dot{\hat x}=A\hat x+B_\ell\tau+L(y-C_j\hat x),\qquad \tau=-K\hat x,\qquad \lambda_{cl}=\lambda(A-B_\ell K)\cup\lambda(A-LC_j).
\end{equation}
Two limits apply to \emph{any} controller:
\begin{itemize}
\item \textbf{Parity interlacing [6]:} for $q_1$ the interval $(2.05,\infty)$ contains two real RHP poles, so a stable controller suffices. For $q_2$, $(2.29,\infty)$ contains one, so the controller itself must be unstable.
\item \textbf{Sensitivity bound [3]:} $\|S\|_\infty\ge\prod_i|z+p_i|/|z-p_i|$, which is $\ge45.9$ for $q_1$ and $\ge187$ for $q_2$, because each zero sits just below the pole $2.24$.
\end{itemize}
In PyBullet the observer-based loop holds $E_3$ only from $1.42^\circ$ (measuring $q_1$) or $0.23^\circ$ (measuring $q_2$), against $4.64^\circ$ with full-state feedback.
\begin{figure}[H]
\centering
\includegraphics[width=\textwidth]{figures/fig_observer.pdf}
\caption{Observer + LQR at $E_3$ in PyBullet using one angle: (a) $\tilde q_1$; (b) torque; (c) largest initial offset held.}
\label{fig:obs}
\end{figure}
\textbf{Answer:} static or PD feedback from one joint angle \textbf{cannot} stabilize the upright, because of the real RHP zero. A dynamic controller \textbf{can}, but only for about a degree with $|\tau|\le20\Nm$, so Phase~3 uses full-state feedback.

% =====================================================================
\section{Phase 3: Control Around the Equilibria}
% =====================================================================
\subsection{Controller Design}
The rank of $\mathcal C$ is 4 at every equilibrium, so all four states can be fed back. Including the torque limit, every controller has the form
\begin{equation}
\tau=\operatorname{sat}(-K\tilde x),\qquad \tilde x=[\tilde q_1,\ \tilde q_2,\ \dot q_1,\ \dot q_2]^T,\qquad |\tau|\le20\Nm .
\end{equation}
The controllers used are:
\begin{itemize}
\item \textbf{LQR} at the upright $E_3$, at $E_2$ and at $E_4$;
\item \textbf{pole placement} at $E_3$, for comparison;
\item \textbf{PD on the elbow angle only} at $E_1$.
\end{itemize}

\subsection{Linear Quadratic Regulator (LQR)}
$K$ minimizes a cost that penalizes both the state error and the torque:
\begin{equation}
J=\int_0^\infty\big(\tilde x^TQ\tilde x+R\,\tau^2\big)\,dt,\qquad A^TS+SA-SB_\ell R^{-1}B_\ell^TS+Q=0,\qquad K=R^{-1}B_\ell^TS .
\end{equation}
\textbf{Expected performance of LQR:}
\begin{itemize}
\item \textbf{Stability:} guaranteed for any $Q\succeq0$, $R>0$ when $(A,B_\ell)$ is controllable.
\item \textbf{Robustness:} phase margin $\ge60^\circ$ and gain margin $[\frac12,\infty)$ [7].
\item \textbf{Tuning:} $Q=I$ weighs 1~rad of angle like 1~rad/s of rate. A smaller $R$ makes the controller faster but demands more torque, so $R$ trades speed against saturation.
\end{itemize}

\subsubsection{Choice of $R$}
For each $R$ we simulated the nonlinear saturated loop from a grid of initial offsets and counted how many were recovered (Fig.~\ref{fig:design}).
\begin{itemize}
\item \textbf{$E_3$:} the recovered area stays at 16--20\% of a $\pm60^\circ$ grid for $10^{-3}\le R\le1$, while the fastest pole moves from $-10$ to $-214$~s$^{-1}$. The limit is $\tau_{\max}$, not the tuning. $R=0.03$ is the knee: 93\% of the best area with a fastest pole of $-40$.
\item \textbf{$E_2$:} the region collapses for $R>0.3$ ($99\%\to90\%\to20\%$), so $R=0.1$.
\item \textbf{$E_4$:} the area peaks at $R\approx0.1$--$0.18$, so $R=0.1$.
\end{itemize}
\begin{figure}[H]
\centering
\includegraphics[width=\textwidth]{figures/fig_design.pdf}
\caption{Weight selection ($Q=I$): (a) share of offsets recovered; (b) fastest closed-loop pole. Stars: chosen $R$.}
\label{fig:design}
\end{figure}

\subsection{Pole Placement at $E_3$ (Comparison)}
Placing the poles at $\{-3,-4,-5,-6\}$ gives $K=[-322.9,\,-134.6,\,-138.1,\,-66.0]$. It settles faster from $2^\circ$ ($1.3$~s vs $4.1$~s) and with less torque ($11.3$ vs $17.6\Nm$), but recovers a smaller region ($20\%$ vs $26\%$ of $\pm40^\circ$; $5^\circ$ vs $6^\circ$ along $q_1$). LQR is therefore kept for $E_3$.

\subsection{$E_1$: PD on the Elbow Angle Only}
At $E_1$ the zeros of $G_2$ are imaginary ($\pm j2.29$), so the sign conflict of~\eqref{eq:pd} does not arise and one angle suffices:
\begin{equation}
\tau=-K_p\tilde q_2-K_d\dot q_2,\qquad K_p=-5,\quad K_d=0.4,
\end{equation}
the pair with the best decay rate on a $(K_p,K_d)$ grid ($\max\operatorname{Re}\lambda=-0.52$).
\textbf{Proof of stability (energy shaping, LaSalle [4]):}
\begin{gather}
V_c=E-E^*+\tfrac12K_p\tilde q_2^{\,2},\qquad \dot V_c=-K_d\,\dot q_2^{\,2}\le0,\\
\nabla^2\big(U+\tfrac12K_p\tilde q_2^{\,2}\big)\big|_{E_1}\ \text{has eigenvalues}\ 0.760,\ 28.575>0 .\nonumber
\end{gather}
$V_c$ has a strict minimum at $E_1$ and never increases, and $\dot q_2\equiv0$ only at $E_1$, so $E_1$ is asymptotically stable. Pure damping ($K_p=0$) decays $19\times$ more slowly.

\subsection{Designed Controllers}
\begin{table}[H]
\centering
\caption{Phase 3 controllers and their closed-loop poles}
\label{tab:ctrl}
\small
\resizebox{\textwidth}{!}{\begin{tabular}{|l|l|l|l|}
\hline
Controller & Design & $K=[k_{q_1},k_{q_2},k_{\dot q_1},k_{\dot q_2}]$ & Closed-loop poles\\ \hline
$E_3$ LQR & $Q=I$, $R=0.03$ & $[-503.0,\,-203.9,\,-218.3,\,-101.5]$ & $-39.86,\,-2.25\pm0.10j,\,-1.34$\\
$E_3$ pole placement & $\{-3,-4,-5,-6\}$ & $[-322.9,\,-134.6,\,-138.1,\,-66.0]$ & $-6,\,-5,\,-4,\,-3$\\
$E_2$ LQR & $Q=I$, $R=0.1$ & $[6.60,\,46.49,\,5.23,\,13.24]$ & $-4.26,\,-2.40,\,-1.24\pm4.16j$\\
$E_4$ LQR & $Q=I$, $R=0.1$ & $[-49.31,\,38.91,\,-13.19,\,6.10]$ & $-5.73,\,-3.21,\,-0.96\pm3.08j$\\
$E_1$ elbow PD & $K_p=-5$, $K_d=0.4$ & $[0,\,-5,\,0,\,0.4]$ & $-0.69\pm2.31j,\,-0.52\pm2.14j$\\ \hline
\end{tabular}}
\end{table}

\subsection{Limits with $\tau_{\max}=20\Nm$}
\subsubsection{Guaranteed Region}
Inside $\Omega_c=\{\tilde x^TS\tilde x\le c\}$ the command never saturates if
\begin{equation}
c^*=\frac{\tau_{\max}^2}{KS^{-1}K^T},\qquad \dot V=-\tilde x^T(Q+K^TRK)\tilde x<0\ \text{ on }\Omega_{c^*}.
\end{equation}
At rest this guarantees $0.72^\circ/1.66^\circ$ ($E_3$), $45.6^\circ/19.9^\circ$ ($E_2$) and $20.0^\circ/23.1^\circ$ ($E_4$) along $q_1$/$q_2$. The bound is conservative because short saturation is usually survived.

\subsubsection{Simulated Region of Attraction}
The real limits come from simulating the nonlinear saturated loop from a grid of initial offsets at rest. A start counts as recovered if it ends within $1^\circ$ and $0.05$~rad/s of $q^*$. The fast integrator used for the grids matches PyBullet on 100\% of random points and 96\% of boundary points.
\begin{figure}[H]
\centering
\includegraphics[width=\textwidth]{figures/fig_roa.pdf}
\caption{Region of attraction (green: recovered). Blue ellipse: guaranteed region. Dashed: offsets with the centre of mass over the pivot.}
\label{fig:roa}
\end{figure}
\begin{table}[H]
\centering
\caption{How far from the equilibrium each controller still works}
\label{tab:limits}
\small
\resizebox{\textwidth}{!}{\begin{tabular}{|l|c|c|c|c|c|}
\hline
Controller & Recovered (grid) & $|\tilde q_1(0)|$ alone & $|\tilde q_2(0)|$ alone & Shoulder kick & Settling\\ \hline
$E_3$ LQR & 26\% ($\pm40^\circ$) & $6^\circ$ & $13^\circ$ & $0.20$ rad/s & $4.1$ s\\
$E_3$ pole placement & 20\% ($\pm40^\circ$) & $5^\circ$ & $9^\circ$ & $0.15$ rad/s & $1.3$ s\\
$E_2$ LQR & 99.7\% ($\pm180^\circ$) & $175^\circ$ & $180^\circ$ & $\ge8$ rad/s & $3.0$ s\\
$E_4$ LQR & 20.5\% ($\pm180^\circ$) & $35^\circ$ & $55^\circ$ & $2.1$ rad/s & $3.6$ s\\
$E_1$ elbow PD & 77.8\% ($\pm180^\circ$) & $40^\circ$ & $115^\circ$ & $2.15$ rad/s & $12.8$ s\\ \hline
\end{tabular}}
\end{table}

\textbf{Why the upright region is a thin band.} Linearizing $\tau_g$ at $E_3$:
\begin{equation}
\tau_{g,1}=9.81\,(2.5\,\tilde q_1+\tilde q_2).
\end{equation}
Gravity pushes the passive shoulder in proportion to the horizontal offset of the total centre of mass. The motor can only counter this through the elbow coupling. Along $\tilde q_2=-2.5\,\tilde q_1$ there is no shoulder torque, and the region reaches the grid edge ($\ge43^\circ$). Across that line, $20\Nm$ is used up after about $6^\circ$.

\subsection{Time Responses}
\begin{figure}[H]
\centering
\includegraphics[width=\textwidth]{figures/fig_e3.pdf}
\caption{$E_3$ in PyBullet: LQR is held from $5.5^\circ$ and falls from $7.5^\circ$; a CoM-balanced offset $(-10^\circ,25^\circ)$ is held with only $3.3\Nm$.}
\label{fig:e3}
\end{figure}
\begin{figure}[H]
\centering
\includegraphics[width=\textwidth]{figures/fig_others.pdf}
\caption{(a) $E_2$ and (b) $E_4$ with LQR; (c) $E_1$ with elbow PD vs.\ open loop; (d) $V_c$ decreasing along the $E_1$ run.}
\label{fig:others}
\end{figure}
\textbf{Observations:}
\begin{itemize}
\item \textbf{Settling time:} pole placement settles fastest at $E_3$, but LQR rescues a larger region. $E_1$ is slowest ($12.8$~s) because its slow mode barely moves the elbow.
\item \textbf{Control effort:} at $E_3$ even a $2^\circ$ offset needs $17.6\Nm$, close to the limit, which explains the small region. At $E_4$ the peak is only $4.5\Nm$.
\item \textbf{Steady-state error:} zero for every controller (each converges to $q^*$ within $0.001^\circ$).
\end{itemize}

% =====================================================================
\section{Phase 4: Swing-Up}
% =====================================================================
\subsection{Energy Target}
Hanging has $E=U(E_1)=-24.525$~J and the upright needs $E_{top}=U(E_3)=+24.525$~J. The controller must therefore add $49$~J with a $20\Nm$ elbow motor, and deliver the arm close enough to $E_3$ for the Phase~3 LQR to catch it.

\subsection{Swing-Up Controller}
\textbf{Partial feedback linearization.} With $h=C\dot q-\tau_g$, row~1 of~\eqref{eq:eom} gives $\ddot q_1=-(M_{12}\ddot q_2+h_1)/M_{11}$. Substituting into row~2, the elbow follows any desired acceleration $v$ exactly:
\begin{equation}
\tau=\Big(M_{22}-\frac{M_{21}M_{12}}{M_{11}}\Big)v+h_2-\frac{M_{21}h_1}{M_{11}}\ \Rightarrow\ \ddot q_2=v .
\end{equation}
\textbf{Energy pumping.} The desired elbow acceleration is
\begin{equation}
v=-k_1\tilde q_2-k_2\dot q_2+k_3\,\dot q_1\,(E-E_{top}) .\label{eq:v}
\end{equation}
\begin{itemize}
\item $k_1$, $k_2$: keep the elbow nearly straight (a stiff and damped elbow).
\item $k_3$: pushes in step with $\dot q_1$ while $E<E_{top}$ (adds energy) and against it when $E>E_{top}$ (removes energy), like a child pumping a swing.
\end{itemize}

\subsection{Hybrid Switching}
\begin{equation}
\tau=\operatorname{sat}\begin{cases}-K\tilde x, & \tilde x^TS\tilde x<c_{sw}=5\quad\text{(catch: Phase 3 LQR at $E_3$)}\\ \text{swing-up law above}, & \text{otherwise}\end{cases}
\end{equation}
If the arm leaves the catch region (a missed catch or a push), swing-up resumes automatically.

\subsection{Choice of Gains}
We searched 144 designs ($k_1\in\{20,50,100\}$, $k_2\in\{2,5,10\}$, $k_3\in\{0.3,1,3,10\}$, $c_{sw}\in\{0.5,1,2,5\}$) from $0.05$~rad off hanging; 22 swung up and held. Each was then rerun from 40 random starts. \textbf{$k=(20,2,3)$, $c_{sw}=5$} succeeded from \textbf{39/40}, with a median first catch of $4.0$~s.

\subsection{Results}
\begin{table}[H]
\centering
\caption{Swing-up runs in PyBullet ($|\tau|\le20\Nm$)}
\label{tab:swing}
\small
\resizebox{\textwidth}{!}{\begin{tabular}{|l|c|c|c|c|}
\hline
Start (from hanging) & First catch & Final deviation from $E_3$ & Torque saturated & Result\\ \hline
$(3^\circ,0)$, $k=(20,2,3)$ & $2.2$ s & $<0.001^\circ$ & 5\% & held\\
$(-10^\circ,20^\circ)$, $k=(20,2,3)$ & $9.5$ s & $<0.01^\circ$ & 17\% & held\\
$(3^\circ,0)$, kick $0.8$ rad/s at 8 s & $2.2$ s, re-caught & $<0.001^\circ$ & 5\% & held\\
$(3^\circ,0)$, textbook $k=(50,5,10)$ & $8.4$ s & $<0.001^\circ$ & 7\% & held (this start)\\ \hline
\end{tabular}}
\end{table}
\begin{figure}[H]
\centering
\includegraphics[width=\textwidth]{figures/fig_swingup.pdf}
\caption{Swing-up from $3^\circ$: (a) angles, (b) energy pumped to $E_{top}$, (c) torque. Green: LQR active.}
\label{fig:swingup}
\end{figure}

\subsection{Why It Works, and Why Other Gains Fail}
\begin{table}[H]
\centering
\caption{Success from 40 random starts vs.\ saturation during swing-up}
\label{tab:sat}
\begin{tabular}{|l|c|c|}
\hline
Gains $(k_1,k_2,k_3)$, $c_{sw}$ & Swing-up time at $|\tau|=20$ & Held (of 40)\\ \hline
$(20,2,3)$, 5 (chosen) & 26\% & \textbf{39}\\
$(50,2,3)$, 5 & 38\% & 4\\
$(100,2,3)$, 5 & 45\% & 2\\
$(20,2,0.3)$, 5 & 25\% & 0\\
$(50,5,10)$, 5 (textbook) & 33\% & 25\\
$(50,5,10)$, 1 (textbook) & 31\% & 13\\ \hline
\end{tabular}
\end{table}
\begin{enumerate}
\item \textbf{Saturation breaks the swing-up law.} The law assumes the full linearizing torque is applied. A stiff elbow (large $k_1$) asks for more than $20\Nm$ more often (26\% $\to$ 45\% of the time). While saturated $\ddot q_2\ne v$, the energy is no longer steered precisely, and the arm reaches the top off the catch region. Success drops from 39 to 2.
\item \textbf{Too little pumping.} With $k_3=0.3$ the energy approaches $E_{top}$ too slowly to reach the catch region within 40~s (0/40).
\item \textbf{The catch is the Phase 3 limit.} The LQR recovers only a thin region around $E_3$. The swing-up only has to deliver the arm into that region, and $c_{sw}=5$ was the most reliable threshold tried.
\item \textbf{Retry.} Because swing-up resumes whenever $\tilde x^TS\tilde x\ge c_{sw}$, a missed catch or a push costs time, not the task (Fig.~\ref{fig:kick}a).
\end{enumerate}
\begin{figure}[H]
\centering
\includegraphics[width=\textwidth]{figures/fig_why.pdf}
\caption{(a) Energy, chosen vs.\ textbook gains; (b) distance to the upright and the catch threshold; (c) success rate over $(k_1,k_3)$.}
\label{fig:why}
\end{figure}
\begin{figure}[H]
\centering
\includegraphics[width=0.9\textwidth]{figures/fig_kick.pdf}
\caption{(a) Knocked off at 8~s, swung up and caught again; (b) phase portrait from $(-10^\circ,20^\circ)$, green where LQR is active.}
\label{fig:kick}
\end{figure}

% =====================================================================
\section{Simulations and Code}
% =====================================================================
\subsection{PyBullet Simulator}
The nonlinear model runs in PyBullet with an interactive dashboard (Fig.~\ref{fig:dash}). The user picks an equilibrium, an initial condition (slider, typed value or by dragging the arm) and a controller: open loop, single-angle PD, single-angle observer + LQR, full-state LQR, pole placement, elbow PD, or swing-up $\to$ LQR. The dashboard shows the nonlinear arm next to the linear model, $\tilde q(t)$, $\tau(t)$ and the closed-loop poles in the $s$-plane, all updating live. Files: \texttt{phase2/acrobot\_phase2\_sim.py}, \texttt{phase3/acrobot\_phase3\_sim.py}, \texttt{phase4/acrobot\_phase4\_sim.py} (instructions in each folder's \texttt{README.md}).
\begin{figure}[H]
\centering
\includegraphics[width=\textwidth]{figures/dashboard.png}
\caption{Interactive PyBullet dashboard during a swing-up from hanging.}
\label{fig:dash}
\end{figure}

\subsection{MATLAB Codes}
Every number in this report was reproduced in MATLAB (Control System Toolbox). The files are listed below; the algorithms they implement follow as pseudo-code.
\begin{table}[H]
\centering
\caption{MATLAB files}
\resizebox{\textwidth}{!}{\begin{tabular}{|l|l|l|}
\hline
Phase & File & Purpose\\ \hline
1 & \texttt{Acrobot\_Nonlinear\_Dynamics.m} & nonlinear model $\dot x=f(x,\tau)$ with the $20\Nm$ limit\\
1 & \texttt{Acrobot\_Linearize.m} & $A$, $B_\ell$ at an equilibrium\\
2 & \texttt{Acrobot\_Phase2\_Main.m} & poles, zeros, ranks, verdicts, growth rates, single-angle PD test\\
3 & \texttt{Acrobot\_Phase3\_Design.m} & the controllers (LQR, pole placement, elbow PD)\\
3 & \texttt{Acrobot\_Phase3\_Main.m} & guaranteed regions, nonlinear runs, region-of-attraction maps\\
4 & \texttt{Acrobot\_Phase4\_Controller.m} & hybrid swing-up $\to$ LQR controller\\
4 & \texttt{Acrobot\_Phase4\_Main.m} & swing-up runs and the 40-start robustness test\\ \hline
\end{tabular}}
\end{table}
\input{pseudo/model}
\input{pseudo/phase2}
\input{pseudo/phase3}
\input{pseudo/phase4}

\subsection{Animations}
Each phase has a 2K Manim animation of the analysis and the simulations:
\href{https://drive.google.com/file/d/19THfoXXzbdMtbteCKE4v0VRLpV869wE0/view?usp=sharing}{\textcolor{blue}{Phase 2}},
\href{https://drive.google.com/file/d/1c-7Yy664ZP4TCJIeDbrAmFozHC8frao1/view?usp=sharing}{\textcolor{blue}{Phase 3}},
\href{https://drive.google.com/file/d/1jPQRn3W-bvUyvSKlHvYmX74bfPX9N16P/view?usp=sharing}{\textcolor{blue}{Phase 4}}.

% =====================================================================
\section{Tabulation}
% =====================================================================
\begin{table}[H]
\centering
\caption{Summary of the stability analysis (Phase 2)}
\begin{tabular}{|c|c|c|c|}
\hline
Equilibrium & RHP poles & Energy shape & Stability\\ \hline
$E_1$ down--down & 0 & minimum & marginally (Lyapunov) stable\\
$E_2$ down--up & 1 & saddle & unstable\\
$E_3$ up--up & 2 & maximum & unstable\\
$E_4$ up--down & 1 & saddle & unstable\\ \hline
\end{tabular}
\end{table}
\begin{table}[H]
\centering
\caption{Summary of the controllers (Phases 2--4)}
\small
\begin{tabular}{|l|l|l|}
\hline
Controller & Works? & Limit with $|\tau|\le20\Nm$\\ \hline
PD on one angle at $E_3$ & no (proved) & always unstable\\
Observer + LQR, one angle, at $E_3$ & yes & $1.42^\circ$ ($q_1$), $0.23^\circ$ ($q_2$)\\
Full-state LQR at $E_3$ & yes & $6^\circ$ / $13^\circ$, $\ge43^\circ$ along the CoM line\\
Pole placement at $E_3$ & yes & $5^\circ$ / $9^\circ$\\
LQR at $E_2$ & yes & almost any rest position (99.7\%)\\
LQR at $E_4$ & yes & $35^\circ$ / $55^\circ$\\
Elbow PD at $E_1$ & yes & 78\% of $\pm180^\circ$\\
Swing-up $\to$ LQR from $E_1$ & yes & 39 of 40 random starts\\ \hline
\end{tabular}
\end{table}

% =====================================================================
\section{Conclusion}
% =====================================================================
The acrobot has four zero-torque equilibria. Only the hanging pose $E_1$ is stable, and only in the Lyapunov sense: with no friction it oscillates forever, while $E_2$, $E_4$ (one RHP pole each) and the upright $E_3$ (two RHP poles) are unstable. The transfer functions and the state matrix agree, and PyBullet confirmed the predicted growth rates to within $0.25\%$. A single measured angle cannot stabilize the upright with PD feedback, because the transfer functions have real right-half-plane zeros. An observer-based controller can, but only for about a degree.

With full-state feedback all four equilibria were stabilized. At the upright, the $20\Nm$ limit rather than the tuning restricts the LQR to a thin band around the line where the centre of mass is over the pivot. Finally, an energy-pumping swing-up with partial feedback linearization brings the arm from hanging to the upright, and the Phase~3 LQR catches and holds it, from 39 of 40 random starts. Stiffer gains fail because they drive the motor into saturation, which breaks the linearization the swing-up relies on.

% =====================================================================
\section{Contributions}
% =====================================================================
% Fill in who did what before submitting.
\begin{enumerate}
\item Report compilation -- \underline{\hspace{4cm}}
\item Phase 2 (stability analysis, Sections 2.1--2.7) -- \underline{\hspace{4cm}}
\item Phase 3 (controller design and limits, Sections 3.1--3.7) -- \underline{\hspace{4cm}}
\item Phase 4 (swing-up, Sections 4.1--4.6) -- \underline{\hspace{4cm}}
\item PyBullet simulator and figures -- \underline{\hspace{4cm}}
\item MATLAB codes -- \underline{\hspace{4cm}}
\item Animations (Manim) -- \underline{\hspace{4cm}}
\end{enumerate}

\begin{thebibliography}{9}
\bibitem{spong} M.~W. Spong, ``The swing up control problem for the Acrobot,'' \emph{IEEE Control Systems Magazine}, vol.~15, no.~1, pp.~49--55, 1995.
\bibitem{tedrake} R.~Tedrake, \emph{Underactuated Robotics}, MIT, 2024. Available: https://underactuated.csail.mit.edu
\bibitem{skogestad} S.~Skogestad and I.~Postlethwaite, \emph{Multivariable Feedback Control}, 2nd ed., Wiley, 2005.
\bibitem{khalil} H.~K. Khalil, \emph{Nonlinear Systems}, 3rd ed., Prentice Hall, 2002.
\bibitem{pybullet} E.~Coumans and Y.~Bai, ``PyBullet, a Python module for physics simulation,'' 2016--2024. Available: https://pybullet.org
\bibitem{youla} D.~C. Youla, J.~J. Bongiorno and C.~N. Lu, ``Single-loop feedback-stabilization of linear multivariable dynamical plants,'' \emph{Automatica}, vol.~10, pp.~159--173, 1974.
\bibitem{anderson} B.~D.~O. Anderson and J.~B. Moore, \emph{Optimal Control: Linear Quadratic Methods}, Prentice Hall, 1990.
\end{thebibliography}

\end{document}
